\documentclass[10pt,a4paper]{amsart}
\usepackage[margin=19mm]{geometry}
\usepackage{amsmath,amssymb,mathtools}
\usepackage[scr=rsfs,cal=euler]{mathalfa}
\usepackage{xfrac}
\usepackage[unicode,hidelinks,bookmarks=false]{hyperref}
\newcommand{\ie}{i.e.\ }
\newcommand{\cf}{cf.\ }
\newcommand{\resp}{respectively\ }
\newcommand{\Subsection}{\subsection}
\providecommand{\nobreakdash}{\nobreak\hbox{-}}
\setlength{\emergencystretch}{2em}
\multlinegap0pt
\usepackage{bm}
\usepackage{bbm}
\usepackage{dsfont}
\usepackage{xspace}
\usepackage{cancel}
\usepackage{mathtools}
\usepackage{amsthm}
\makeatletter
\@ifundefined{thm}{\newtheorem{thm}{Thm}}{}
\@ifundefined{lem}{\newtheorem{lem}[thm]{Lemma}}{}
\@ifundefined{prop}{\newtheorem{prop}[thm]{Proposition}}{}
\makeatother
\newcommand{\eps}{\varepsilon}
\title[Many Antipodal Pairs Force Many Neighboring Pairs]{Many Antipodal Pairs Force Many Neighboring Pairs}
\author{G\'abor Dam\'asdi, Laurentiu Ploscaru}
\date{Source: arXiv:2608.02605v1 (CC BY 4.0)}
\begin{document}
\maketitle

\noindent\emph{Source and licence.} Many Antipodal Pairs Force Many Neighboring Pairs. Version arXiv:2608.02605v1 (CC BY 4.0), distributed under the Creative Commons Attribution 4.0 International licence, as recorded in the arXiv licence metadata for this exact version and shown on the abstract page https://arxiv.org/abs/2608.02605v1. Full rights notice: licenses/arxiv-2608.02605-CC-BY-4.0.txt.\par\smallskip

\noindent\textit{Editorial scope.} Counted unit: the Prop whose source label is \texttt{ref:annularconstr} in the article (source0.tex lines 557--563, source label \texttt{ref:annularconstr}), delivered together with the article's own complete original proof of it (source0.tex lines 590--623, one \texttt{proof} environment of 34 lines). No mathematics is written, repaired or re-derived: statement and proof are the article's own text line for line. The proof is detached from the statement in the source: the article states the result at lines 557--563 and prints its proof at lines 590--623, and the proof environment's own header names the statement, so the association is the article's own. Required context retained uncounted: lem:annulus-count (lines 578--581). Every definition, display and label of those lines is retained unchanged. Omitted from this package and not redistributed: 1--556, 564--577, 582--589, 624--757. Nothing inside a retained excerpt is omitted or altered: every retained body is delivered exactly as the article's lines are written, so no omission receipt of any kind accompanies this package. Citations: the retained excerpts carry no citation command at all, so no bibliography entry of the article is transcribed and no reference is redistributed. Reference adaptation: none was needed and none was made; every reference inside the delivered unit resolves inside it, so no unresolved reference or citation is produced. Printed slips kept literal: no printed word, symbol, hypothesis or number of the counted statement or of its proof was altered. Layout and class: the article's own class is \texttt{article} with packages including amsfonts, amsmath, amsthm, bm, bbm, thmtools, thm-restate, latexsym, color, epsfig, mathrsfs, enumerate, appendix, fancyhdr; the wrapper is the lane's pinned \texttt{amsart} editorial wrapper at 10pt on A4, which restates the theorem-like environments the retained text needs and adds the article's own macro definitions that the retained text needs (\texttt{\textbackslash eps}), with extra wrapper packages (bm, bbm, dsfont, xspace, cancel, mathtools). The delivered numbering is the unit's own. Assets: no retained span contains a figure environment, a graphics-inclusion command, a PDF-page inclusion, a table or a TikZ picture; no image file is copied into this package, the package declares no asset dependency and no image file is redistributed with it. Licence: version arXiv:2608.02605v1 of this article is distributed under the Creative Commons Attribution 4.0 International licence, as recorded in the arXiv licence metadata for this exact version and shown on the abstract page https://arxiv.org/abs/2608.02605v1. A full rights notice is delivered at licenses/arxiv-2608.02605-CC-BY-4.0.txt. Characters: every retained excerpt is pure ASCII, so the delivered engine, XeLaTeX with its default Latin Modern text font, reports no missing character.
\begin{lem}\label{lem:annulus-count}
        Let $-1\leq \alpha_1<\alpha_2<1$, $0\leq \delta_1\leq \delta_2\leq \delta<1$ and $r,s\geq 2$ be integers. Then: 
        $$\big|AU(\delta,r,s)\cap AS(\alpha_1,\alpha_2,\delta_1,\delta_2)\big|=\big(1+\lfloor \alpha_2r\rfloor-\lceil \alpha_1r\rceil\big)\cdot\big(1+\lfloor \delta_2s/\delta\rfloor-\lceil  \delta_1s/\delta\rceil\big).$$
\end{lem}

\begin{prop}\label{ref:annularconstr}
    Let $0<2\eps<\delta\leq 1/2$ and let $r,s\geq \eps^{-2}$ be integers. Then:
    \begin{align*}
     \left|\mathcal{A}_{\delta}\big(AU(\delta,r,s)\big)\right|=&\ \Theta\left(\sqrt{\delta}\cdot r^2s^2\right); \\
     \left|\mathcal{N}_{\eps}\big(AU(\delta,r,s)\big)\right|=&\ \Theta\left((\eps\cdot rs)^2/\delta\right).
\end{align*}
\end{prop}

\begin{proof}[Proof of Proposition \ref{ref:annularconstr}]
Set $X:= AU(\delta,r,s)$ and start with estimating $\mathcal{N}_\eps(X)$. We claim that for every point $Q\in X$ the ball of radius $\eps$ centered at $Q$ contains $\Theta\left(\eps^2rs/\delta\right)$ other points of $X$. Since each of these points forms an $\eps$-neighboring pair with $Q$, and $|X|=2rs$, the claim does imply that $|\mathcal{N}_\eps(X)|=\Theta\left(\eps^2r^2s^2/\delta\right)$.

Due to the circular symmetry of our construction, we can assume the point $Q$ has coordinates $\big((s-\delta q)/2s,0\big)$ for some $q\in \mathbb{Z}_s$. Let $(k,l)\in  \mathbb{Z}[-r,r-1] \times \mathbb{Z}_s$ and let $P\in AU(\delta,r,s)$ be a point with coordinates
$(1/2-\delta l/2s)\big(\cos(\pi k/r),\sin(\pi k/r)\big)$. Then the distance $|PQ|$ satisfies: 
    \begin{align*}
        |PQ|^2 &= \left(\dfrac{1}{2}-\dfrac{\delta q}{2s}-\cos\left(\dfrac{\pi k}{r}\right)\cdot \left(\dfrac{1}{2}-\dfrac{\delta l}{2s}\right)\right)^2 +\left(\dfrac{1}{2}-\dfrac{\delta l}{2s}\right)^2\cdot \sin^2\left(\dfrac{\pi k}{r}\right) \\
         &=  \left(\dfrac{1}{2}-\dfrac{\delta q}{2s}\right)^2 - 2\cos\left(\dfrac{\pi k}{r}\right)\cdot \left(\dfrac{1}{2}-\dfrac{\delta q}{2s}\right)\left(\dfrac{1}{2}-\dfrac{\delta l}{2s}\right) + \left(\dfrac{1}{2}-\dfrac{\delta l}{2s}\right)^2\\
         &= \dfrac{\delta^2}{4s^2}\cdot (q-l)^2 + 2\left(1-\cos\left(\dfrac{\pi k}{r}\right)\right)\cdot \left(\dfrac{1}{2}-\dfrac{\delta q}{2s}\right)\left(\dfrac{1}{2}-\dfrac{\delta l}{2s}\right).
    \end{align*}
   \indent Let us note that $\dfrac{1}{2}\geq 2\left(\dfrac{1}{2}-\dfrac{\delta q}{2s}\right)\left(\dfrac{1}{2}-\dfrac{\delta l}{2s}\right)\geq\dfrac{1}{2}(1-\delta)^2=\dfrac{1}{8}$ since $1-\delta\geq \dfrac{1}{2}$. Moreover, we know that $x^2/4\leq 1-\cos x\leq x^2/2$ whenever $|x|<2.7$, which will be useful later as well. 
 
 Consequently, if $|k|\geq 2\varepsilon r$, then $|PQ|^2\geq (1-\cos(\pi k/r))/8\geq 2(\pi k/8r)^2>\varepsilon^2$. Similarly, if $|q-l|>2\varepsilon s/\delta$, then $|PQ|^2>\varepsilon^2$. This means that the ball of radius $\eps$ centered at $Q$ is contained in $AS(-2\eps r,\ 2\eps r,\ \delta q/s-\eps,\ \delta q/s +\eps)$. 
 
  On the other hand, if both $|q-l|\leq 6\eps s/5\delta$ and $|k|\leq \eps r/2$ hold, then we deduce from the above that $|PQ|^2\leq 9\eps^2/25 +\pi^2k^2/(2r)^2\leq 9\eps^2/25+16\eps^2/25=\eps^2$, which means that the ball of radius $\eps$ centered at $Q$ contains $AS(-\eps r/2,\ \eps r/2,\ \delta q/s-3\eps/5,\ \delta q/s+3\eps/5)$. It follows by Lemma \ref{lem:annulus-count} that any point $Q\in X$ contributes with $\Theta\left(\eps^2rs/\delta\right)$ pairs to $\mathcal{N}_{\eps}(X)$.

We now move to $\mathcal{A}_\delta(X)$. Similarly as before, we claim that for every point $Q\in X$ there are $\Theta\left(\sqrt{\delta}\cdot rs\right)$ points in $X$ which are $\delta$-antipodal with $Q$. Again, as $|X|=2rs$, this claim implies that $|\mathcal{A}_\delta(X)|=\Theta\left(\sqrt{\delta}\cdot r^2s^2\right)$.

We may assume this time that $Q$ has coordinates $\big((\delta q-s)/2s,0\big)$ for some $q\in \mathbb{Z}_s$, i.e. it is situated on the negative part of the $x$-axis due to the symmetry. We let $(k,l)\in \mathbb{Z}[-r,r-1]\times \mathbb{Z}_s$ and let $P\in AU(\delta,r,s)$ be the point with coordinates
$(1/2-\delta l/2s)\big(\cos(\pi k/r),\sin(\pi k/r)\big)$. We estimate $|PQ|$ in a similar fashion: 
 \begin{align*}
        |PQ|^2 &= \left(\dfrac{1}{2}-\dfrac{\delta q}{2s}+\cos\left(\dfrac{\pi k}{r}\right)\cdot \left(\dfrac{1}{2}-\dfrac{\delta l}{2s}\right)\right)^2 +\left(\dfrac{1}{2}-\dfrac{\delta l}{2s}\right)^2\cdot \sin^2\left(\dfrac{\pi k}{r}\right) \\
         &=  \left(\dfrac{1}{2}-\dfrac{\delta q}{2s}\right)^2 + 2\cos\left(\dfrac{\pi k}{r}\right)\cdot \left(\dfrac{1}{2}-\dfrac{\delta q}{2s}\right)\left(\dfrac{1}{2}-\dfrac{\delta l}{2s}\right) + \left(\dfrac{1}{2}-\dfrac{\delta l}{2s}\right)^2\\
         &= \left(1-\delta\cdot \dfrac{q+l}{2s}\right)^2  - 2\left(1-\cos\left(\dfrac{\pi k}{r}\right)\right)\cdot \left(\dfrac{1}{2}-\dfrac{\delta q}{2s}\right)\left(\dfrac{1}{2}-\dfrac{\delta l}{2s}\right).
    \end{align*}
  \indent Recalling the useful bounds from above, we deduce that $|PQ|^2\leq 1-\big(1-\cos(\pi k/r)\big)/8\leq 1-\pi^2k^2/32r^2$. Therefore, if $|k|\geq 8r\sqrt{\delta}/\pi$, then $|PQ|^2< 1-2\delta<(1-\delta)^2$, which implies that all points of $X$ which form $\delta$-antipodal pairs with $Q$ must be in $AS(-8r\sqrt{\delta}/\pi,8r\sqrt{\delta}/\pi,0,\delta)$.  

 Finally, we claim that if $2l\leq s$ and $|k|\leq \sqrt{2\delta}\cdot r/\pi$, then $|PQ|^2\geq (1-\delta)^2$, which would imply that all points of $X$ from inside $AS(-\sqrt{2\delta}\cdot r/\pi,\sqrt{2\delta}\cdot r/\pi,0,\delta/2)$ are $\delta$-antipodal with $Q$. Consequently, Lemma \ref{lem:annulus-count} applied to these two annulus sectors would then give us that there are $\Theta\left(\sqrt{\delta}\cdot rs\right)$ such points, as desired.

 To prove this last claim, we have $\left(1-\delta\cdot \dfrac{q+l}{2s}\right)^2\ge 4\left(\dfrac{1}{2}-\dfrac{\delta q}{2s}\right)\left(\dfrac{1}{2}-\dfrac{\delta l}{2s}\right)$ by AM-GM,
 while also $2\big(1-\cos(\pi k/r)\big)\leq \pi^2k^2/r^2$. It follows that
 $$|PQ|^2\geq \left(4-\dfrac{\pi^2k^2}{r^2}\right)\left(\dfrac{1}{2}-\dfrac{\delta q}{2s}\right)\left(\dfrac{1}{2}-\dfrac{\delta l}{2s}\right)\geq \left(1-\dfrac{\pi^2k^2}{4r^2}\right)\left(1-\dfrac{\delta}{2}\right)(1-\delta).$$
  Therefore, it is enough to prove that $1-(\pi k/2r)^2\geq (1-\delta)(1-\delta/2)^{-1}=1-\delta(2-\delta)^{-1}$. But this is true since $(\pi k/2r)^2\leq (\sqrt{2\delta}/2)^2=\delta/2<\delta(2-\delta)^{-1}$, finishing the proof.
\end{proof}

\end{document}
